% !TEX root = ../../main.tex
\subsection{系统总体方案}

FoodShield 是一个面向外卖订单通信场景的隐私认证与可信审计系统。系统以外卖平台中用户与骑手之间的订单沟通为基本应用场景，围绕身份信息暴露、订单隐私泄露、聊天记录篡改和纠纷追责困难等问题，设计了集匿名身份、订单认证、加密通信、日志完整性验证和条件溯源于一体的安全通信方案。

传统外卖订单通信系统通常更关注业务连通性，即用户和骑手能否完成消息交互，而对通信过程中的安全边界关注不足。例如，骑手端可能接触到超出配送所需范围的用户信息，普通订单号可能被误用或伪造，管理员端可能拥有过大的信息查看权限，历史通信记录也可能在纠纷处理前被修改、删除或重排。针对上述问题，FoodShield 将“订单”作为系统安全控制的核心索引，将用户身份、订单资格、通信通道和审计日志统一绑定到订单生命周期中，从而在业务可用性的基础上增加隐私保护和可信审计能力。

系统总体设计遵循“日常匿名、订单认证、日志可信、异常可溯源” 的原则。在日常通信过程中，系统不向骑手端暴露用户真实身份，也不向骑手端展示用户 PID；用户注册阶段不再是简单输入用户名生成匿名标识，而是采用“用户名、口令、机器识别”相结合的方式，降低机器批量注册和匿名身份滥用风险。用户创建订单后，系统将订单号、用户匿名身份、骑手身份和订单状态进行绑定，并生成基于 HMAC-SM3 的订单 Token。用户进入订单通信通道前，需要提交订单号和 Token，服务器验证通过后才允许其加入对应订单房间。

通信过程中，用户端与骑手端通过订单级 WebSocket 通道进行实时通信。系统对通信消息进行加密存储，并对订单号、发送方角色、时间戳、消息序号和密文摘要等字段计算消息哈希。审计阶段，系统将消息摘要作为 Merkle Tree 叶子节点，生成 Merkle Root 并通过快照机制保存。后续若通信记录被修改、删除、插入或重排，重新计算得到的 Root 将与历史快照不一致，从而能够检测日志篡改行为。

当发生投诉、恶意骚扰、虚假争议或平台审计需求时，系统支持用户端和骑手端基于订单号发起条件溯源申请。管理员登录后只能默认查看订单摘要、消息哈希、Merkle Root 和审计记录，不直接展示通信明文内容；只有在满足权限校验、审计理由和操作记录要求后，管理员才能基于订单号触发条件溯源。该设计避免了直接通过 PID 任意查询真实身份的问题，使系统在保护用户隐私的同时保留必要的平台管理能力。

从功能组成上看，FoodShield 主要包括以下核心模块：

\begin{enumerate}[label=\arabic*., leftmargin=2em]
\item \textbf{用户注册与身份初始化模块}：负责用户注册、口令认证、机器识别、PID 生成和匿名身份初始化，避免系统仅具有标识而缺少真实认证流程。
\item \textbf{匿名身份管理模块}：负责生成和维护用户 PID，并保存 PID 与真实身份之间的受控映射关系。PID 默认不在普通前端页面展示，骑手端不可见。
\item \textbf{订单 Token 认证模块}：负责根据订单号、PID、时间戳和随机数生成订单认证 Token，并在用户进入通信通道前进行合法性、归属关系和时效性验证。
\item \textbf{订单级实时通信模块}：负责用户端与骑手端之间的 WebSocket 通信、订单房间管理、消息转发、订单历史记录和订单搜索。
\item \textbf{消息保护与日志审计模块}：负责消息加密存储、SM3 摘要生成、Merkle Tree 构建、Merkle Root 快照保存和日志完整性验证。
\item \textbf{条件溯源与管理员审计模块}：负责投诉申请、管理员权限校验、基于订单号的条件溯源和管理员操作审计记录。
\end{enumerate}

在技术实现上，系统采用 Python Flask 作为后端服务框架，使用 SQLite 存储用户、订单、消息和审计数据，使用 WebSocket 实现实时通信。在密码学机制方面，系统使用 HMAC-SM3 构造订单认证 Token，使用 SM4 对通信消息进行加密存储，使用 SM3 生成消息摘要，并使用 Merkle Tree 对通信日志进行完整性验证。上述机制共同构成了 FoodShield 的安全基础，使其能够在外卖订单通信这一具体场景中兼顾隐私保护、业务可用性和平台管理能力。
